\subsection[The \boldmath NNLO FTapprox combination]{The \boldmath $\textup{NNLO}_{\textup{FTapprox}}$ combination}
%\subsection{The NNLO\texorpdfstring{$_{\mbox{FTapprox}}$}{FTapprox} combination}
\label{sec:NNLO_FTapprox}

%\textit{Authors:} Massimiliano Grazzini, Gudrun Heinrich, Stephen Jones, Stefan Kallweit, Matthias Kerner, Jonas M. Lindert, Javier Mazzitelli (1803.02463)\\

\textbf{Massimiliano Grazzini, Gudrun Heinrich, Stephen Jones, Stefan Kallweit, Mat\-thias Kerner, Jonas M. Lindert, Javier Mazzitelli~\cite{Grazzini:2018bsd}.}



%summarize the approach and related uncertainties of Ref.~\cite{Grazzini:2018bsd}. (GH, MG, ToDo)


\newcommand{\nnloBP}{NNLO$_{\mathrm{B-proj}}$}
\newcommand{\nnloNI}{NNLO$_{\mathrm{NLO-i}}$}
\newcommand{\nnloFT}{NNLO$_{\mathrm{FTapprox}}$}
\newcommand{\nloFT}{$\mathrm{NLO}_{\mathrm{FTapprox}}$}

\noindent QCD corrections to $HH$ production beyond NLO have so far only been computed in approximate form since the relevant two- and three-loop amplitudes for a full NNLO calculation are not yet available. In Ref.~\cite{Grazzini:2018bsd} an approximate NNLO result, dubbed \nnloFT, was presented. In this approach, the exact NLO contribution of Ref.~\cite{Borowka:2016ehy} is incorporated. The approximation of the NNLO correction was constructed starting from the observation that the double-real emission contribution to the NNLO cross section requires only one-loop amplitudes, and can thus be computed exactly with full top quark mass dependence, for example, by using OpenLoops \cite{Cascioli:2011va,Buccioni:2019sur}. However, the approximation of the higher-loop amplitudes entering the NNLO corrections must be done with care, in order not to jeopardize the cancellation of IR singularities.

The approximation is defined as follows. Working in the HTL for each $n$-loop squared amplitude that needs to be computed for a given partonic subprocess, ${\cal A}^{(n)}_{\rm HTL}$, the reweighting factor ${\cal A}_{\rm Full}^{\rm Born}/{\cal A}_{\rm HTL}^{\rm (0)}$ is applied. Here ${\cal A}_{\rm Full}^{\rm Born}$ stands for the exact lowest order (loop-induced) squared amplitude for the corresponding subprocess. The NNLO corrections computed in this way are expected to provide a better estimate of the full result compared to estimates based on a naive reweighting procedure. The impact of NNLO QCD corrections computed in this way ranges from $+11\%$ to $+7\%$ as the collider energy ranges from $13$ to $100$\,TeV \cite{Grazzini:2018bsd}.

Once the employed approximation is defined, the important next question is the size of the remaining uncertainty due to finite top quark mass effects that are still missing compared to an NNLO (3-loop) calculation with full top quark mass dependence.

As a first step to define an uncertainty, we can check the quality of the analogous \nloFT\ approximation, where one can compare to an exact result. At $14$\,TeV the difference between \nloFT\ and the exact NLO cross section is about $4\%$, or, considering only the higher-order corrections, the relative difference between the \nloFT\ corrections and their exact calculation is about $11\%$, which can be translated into a lower bound on the uncertainty of the NNLO cross section. Considering the smaller impact of NNLO with respect to NLO corrections, this translates into a $1.2\%$ effect at NNLO. To be conservative, one can multiply this estimate by a factor of two, thereby obtaining an uncertainty that ranges from $\pm 2.3\%$ at $\sqrt{s}=13$\,TeV to $\pm 3.1\%$ at $\sqrt{s}=100$\,TeV.
However, the study of Ref.~\cite{Grazzini:2018bsd} shows that the difference between the ``best'' prediction, \nnloFT, and the prediction obtained through a simple reweighting of the $HH$ invariant-mass distribution, ${\rm NNLO}_{\rm NLO-I}$ \cite{Borowka:2016ypz}, increases with the energy. Therefore, the uncertainty of the \nnloFT\ prediction is finally defined by taking the half difference of the reference \nnloFT\ result with the ${\rm NNLO}_{\rm NLO-I}$ result. This leads to a final uncertainty due to missing finite top quark mass effects that ranges from $\pm 2.6\%$ at $\sqrt{s}=13$ TeV to $\pm 4.6\%$ at $\sqrt{s}=100$ TeV, and it is therefore slightly more conservative, especially at very high centre-of-mass energies.

Due to the inclusion of the maximal amount of top quark mass dependence in all ingredients, i.e.\ except for those virtual diagrams which are not yet available, the \nnloFT\ prediction had been the recommendation of the LHC Higgs Working Group up to 2026.
